\subsection{The ninth-root base and its exact class norm quotient}

\begin{statement}{A.191}{Proved}
{A cyclotomic base change with a nonzero localized class quotient}
Put \(d=649\), choose the positive real roots
\(\alpha^3=3\), \(t^3=d^2\alpha\), and fix a primitive
ninth root \(z\) with \(\zeta=z^3\).  Define
\[
 F=\Q(\alpha,\zeta),\quad F'=\Q(\alpha,z),\quad
 E_2=\Q(t,\zeta),\quad M=\Q(t,z).
\tag{A.191.1}
\]
These are the indicated embedded fields of the fixed tower.
In particular \([F':\Q]=18\), \([M:\Q]=54\).  The two
cyclic cubic generators over \(F\) are distinguished by
\[
 \sigma(t)=\zeta t,\quad \sigma(z)=z,
 \qquad h(t)=t,\quad h(z)=z^4.
\tag{A.191.2}
\]
Thus \(\langle\sigma\rangle=\operatorname{Gal}(M/F')\),
\(\langle h\rangle=\operatorname{Gal}(M/E_2)\), and they commute.
Let \(S\) consist of the primes above \(2,3,11,59\) and
the archimedean places, and put
\[
 A=\operatorname{Cl}(M)\{3\},\quad
 A_S=\operatorname{Cl}_S(M)\{3\},\quad
 A_2=\operatorname{Cl}(E_2)\{3\}.
\tag{A.191.3}
\]
The actual unit field-norm index and invariant class order are
\[
 [\mathcal O_{F'}^\times:
    \mathcal O_{F'}^\times\cap N_{M/F'}M^\times]=27,
 \qquad |A^\sigma|=27.
\tag{A.191.4}
\]
The ideal norm \(N_{M/E_2}\), with the residue-degree weights
in its usual definition, factors through localization and induces
the following isomorphisms:
\[
 \frac{A}{(\sigma-1)A}
   \xrightarrow[\mathrm{localization}]{\ \sim\ }
 \frac{A_S}{(\sigma-1)A_S}
   \xrightarrow[\ N_{M/E_2}\ ]{\ \sim\ }
 A_2\simeq(\Z/3)^3.
\tag{A.191.5}
\]
Ideal extension \(A_2\to A_S\) is zero.
Both \(A\) and \(A_S\) are annihilated by
\(1+\sigma+\sigma^2\); as finite modules over
\(\Z_3[\zeta]\), they require exactly three generators modulo
the uniformizer \(\sigma-1\).  The asserted maps determine
these coinvariants, without specifying the lengths of the three
cyclic module factors.
\end{statement}

\begin{proof}
\emph{Fields and the prime above three.}
The element \(t\) satisfies \(t^9=3d^6\).
A cube root of \(3d^6\) in an abelian extension of \(\Q_3\)
would contain \(\Q_3(\alpha)\), an Eisenstein cubic extension
which is not Galois.  This is impossible for a subfield of an
abelian extension.  Kummer theory over \(\Q_3(z)\) therefore
gives radical degree nine.  Its cyclotomic degree is six.
These equal the global degrees, so \(M\) has full local
decomposition group at three, namely
\((\Z/9)\rtimes(\Z/9)^\times\).
Conjugation by the unit \(-1\) puts the whole translation
subgroup in the commutator.  The maximal abelian subfield is
therefore the cyclotomic field.  Since that field is totally
ramified at three, the unramified quotient is trivial.
Thus \(M/\Q_3\) is totally ramified.  In particular the
cyclic extension \(M/E_2\) is ramified there.
The same maps prove (A.191.1)--(A.191.2) and show that
\(F'/\Q(\zeta)\) is a Galois three-extension ramified only
above three.  The inertia-normal-generation argument of
A.138 gives
\[
                 \operatorname{Cl}(F')\{3\}=0.
\tag{A.191.6}
\]
Indeed the three-class Hilbert extension, followed by its
normal closure over the class-number-one field \(\Q(\zeta)\),
is a finite three-extension ramified only there.  Its inertia
normally generates, hence equals the group by the maximal-subgroup
property of finite three-groups.  An extension unramified over
\(F'\) must consequently be trivial.

\emph{Literal units and their exact character matrix.}
Here are small definitions of the nine units used in the index
calculation.  For a six-tuple \(c\), write
\(c(z)=\sum_{j=0}^5c_jz^j\).  Define
\(e_i=(a_i(z)+\alpha b_i(z)+\alpha^2c_i(z))/3\)
by the table; its first row is \(e_1=z\).
\begingroup
\setlength{\arraycolsep}{3pt}
\[
\begin{array}{c|r r r r r r|r r r r r r|r r r r r r}
 i&\multicolumn{6}{c|}{a_i}&\multicolumn{6}{c|}{b_i}&\multicolumn{6}{c}{c_i}\\ \hline
1&0&3&0&0&0&0&0&0&0&0&0&0&0&0&0&0&0&0\\
2&0&0&0&0&0&0&0&0&0&0&0&0&1&0&1&1&-1&1\\
3&0&0&0&0&0&0&0&0&0&0&0&0&0&-1&1&1&-1&0\\
4&1&0&-1&-1&0&-2&1&-1&0&0&-1&1&0&0&0&0&-1&0\\
5&-1&2&2&-2&1&1&0&2&0&0&1&0&0&0&0&0&0&0\\
6&0&-1&-2&0&-2&-1&1&1&1&0&0&0&-1&0&0&-1&1&0\\
7&-2&2&0&-1&1&0&0&-1&0&-1&0&-1&0&-1&0&0&-1&0\\
8&-2&0&-1&-1&0&-2&-1&1&-1&-1&0&-1&-1&0&0&0&0&0\\
9&1&2&0&2&1&0&1&-1&1&1&0&1&0&0&-1&0&0&-1
\end{array}
\tag{A.191.7}
\]
\endgroup
All arithmetic in this table takes place in the explicit algebra
\(\Q[\alpha,z]/(\alpha^3-3,z^6+z^3+1)\), which is the
field \(F'\).  The multiplication characteristic polynomial of
each \(e_i\) is monic integral with constant term one.  Thus
each element and its inverse are integral.  The retained rational
certificate checks these identities, and also the following
finite saturation test without a unit-group completeness input.

At a residue map \((\alpha,z)\mapsto(a,b)\) to \(\F_p\),
record the row whose \(i\)-th entry is defined by
\(e_i(a,b)^{(p-1)/3}=b^{3r_i}\).
The nine maps
\[
\begin{array}{c|r r}
p&a&b\\ \hline
271&117&106\\271&24&106\\271&130&248\\
271&117&258\\271&24&125\\271&24&169\\
271&130&169\\307&192&287\\577&329&546
\end{array}
\quad\text{give}\quad
\begin{pmatrix}
0&1&2&0&2&2&1&1&1\\
0&1&2&1&0&1&1&1&0\\
0&2&1&1&1&0&1&1&2\\
0&2&0&0&1&0&2&2&2\\
0&1&0&0&1&2&2&2&2\\
0&0&2&2&0&2&2&0&0\\
0&0&2&0&2&1&1&2&2\\
1&2&2&2&1&1&2&0&0\\
1&1&0&1&2&0&1&1&2
\end{pmatrix}.
\tag{A.191.8}
\]
The matrix is invertible over \(\F_3\).  The free unit rank
of \(F'\) is eight, and its three-primary roots of unity are
exactly \(\mu_9\).  For the latter assertion the abelianization
of \(\operatorname{Gal}(F'/\Q)=(\Z/3)\rtimes(\Z/9)^\times\)
is \((\Z/9)^\times\), so it contains no larger cyclotomic
three-power field.  Thus the unit quotient modulo cubes has
dimension nine.  The displayed rows prove that the \(e_i\)
span it and generate the pro-three unit completion.

For each \(q=11,59\), the field \(F'\) has three primes
with residue field \(\F_{q^6}\).  The radicand
\(t^3=d^2\alpha\) has valuation two at each, so all six
primes are totally ramified in \(M/F'\).  Together with the
prime at three these are exactly its seven ramified primes.
The cube-character rows of the \(e_i\), with target generator
the residue of \(\zeta=z^3\), are, in one ordering of the primes,
\[
\begin{array}{c|rrrrrrrrr}
11&1&0&1&1&2&1&1&1&2\\
11&1&0&1&1&1&2&2&0&2\\
11&1&0&1&1&1&2&1&0&2\\
59&2&0&2&2&1&2&1&2&1\\
59&2&0&2&2&1&2&0&2&1\\
59&2&0&2&2&0&0&1&1&1
\end{array}
\tag{A.191.9}
\]
This matrix has rank three.  At the six tame ramified primes,
a unit is a field norm exactly when its row value is zero.
Units are norms at every other finite prime except possibly
three; at three its character is determined by global
reciprocity.  The cyclic Hasse norm theorem, in the precise form
used in A.175, therefore gives the index in (A.191.4).
The ambiguous-class calculation (A.173.7)--(A.173.10) applies
with this base field and seven ramified primes.  Its unit
Herbrand quotient remains \(1/3\): the rational logarithmic
unit representation is \(\Q[\langle\sigma\rangle]^9-\Q\).
Using (A.191.6) and the index \(3^3\), it gives
\[
                       |A^\sigma|=3^6/3^3=27.
\tag{A.191.10}
\]

\emph{The ideal norm and the localization map.}
Let \(H/E_2\) be the full Hilbert three-class field.
A.173 gives \(\operatorname{Gal}(H/E_2)\simeq(\Z/3)^3\),
with the three explicit unramified characters
\(\chi_0,\chi_1,\chi_2\) of A.170 and A.173.
The latter two have radicals
\(11(1+6\alpha)\) and \(649\).
Because \(M/E_2\) is cyclic cubic ramified at three,
\(H\cap M=E_2\).  Consequently \(HM/M\) is unramified
of degree 27.  It splits at every prime above \(S\):
\(\chi_0\) splits at two by A.170, and the other two
radicals reduce to one there, with cubing an automorphism
on principal two-units.  At three all three radicals are
literal local cubes by A.170 and A.173.  At either changed
prime \(q\), the extension \(M/E_2\) is unramified of
residue degree three, since \(q^2\) has exact order three
modulo nine.  Its Frobenius character is therefore three
times the old character and vanishes.

Here is the corresponding ideal map without a class-field
identification left implicit.  The norm of a prime above a
changed prime is three times its old prime class.  The target
\(A_2\) has exponent three.  Every old two- or three-prime
class is zero, since all three detecting characters vanish
there.  Thus the norm kills the subgroup generated by the
\(S\)-prime classes and factors as
\[
                A\longrightarrow A_S
                       \xrightarrow{\ N_{M/E_2}\ } A_2.
\tag{A.191.11}
\]
It is onto: for each element of \(\operatorname{Gal}(H/E_2)\),
Chebotarev in the disjoint compositum \(HM/E_2\) supplies
a prime split in \(M\) with that Hilbert Frobenius.
A degree-one prime above it has exactly that ideal norm.
Taking primary parts retains surjectivity.  Conversely,
\(A_2\) is generated by its changed-prime classes, so their
extended ideals are supported on \(S\) and the extension
map \(A_2\to A_S\) is zero.

The norm is \(\sigma\)-equivariant, and its target action
is trivial by A.173.  For the finite group \(A\), the
kernel and cokernel of \(\sigma-1\) have equal order.
By (A.191.10) its coinvariant group has order 27.
The surjection (A.191.11) therefore proves both
isomorphisms in (A.191.5).  Finally
\(1+\sigma+\sigma^2=j_{M/F'}N_{M/F'}\) on ideal classes,
and (A.191.6) kills its three-primary image.  This proves
the module assertions, including the three-dimensional
quotient by \(\sigma-1\).

The certificate
\path{code/two_inert_zeta9_base_units_certificate.py}
and its retained output/data verify the literal field and unit
identities, the full nine-column saturation, and all six tame
rows.  The small table (A.191.7) is checked against these
identities in the natural \((\alpha,z)\) basis.  The proof
uses only the indicated invariant and coinvariant class maps;
it supplies no further unit-character value from their orders.
\end{proof}
